%% ============================================================
%%  ES-RL Technical Document
%%  Evolution Strategies with RL-style Gradient Estimation
%%  Compile in Overleaf (pdflatex, no bibtex needed — refs inline)
%% ============================================================

\documentclass[12pt, a4paper]{article}

%% ── Packages ────────────────────────────────────────────────
\usepackage[T1]{fontenc}
\usepackage[utf8]{inputenc}
\usepackage{lmodern}
\usepackage[margin=2.5cm, headheight=15pt]{geometry}

%% Math
\usepackage{amsmath, amssymb, amsthm, mathtools}

%% Floats & tables
\usepackage{graphicx}
\usepackage{booktabs}
\usepackage{array}

%% Algorithms
\usepackage{algorithm}
\usepackage{algpseudocode}
\renewcommand{\algorithmicrequire}{\textbf{Require:}}

%% Colours & boxes
\usepackage[dvipsnames]{xcolor}
\usepackage{tcolorbox}
\tcbuselibrary{theorems, skins, breakable}

%% Hyperlinks
\usepackage[colorlinks=true,
            linkcolor=MidnightBlue,
            citecolor=ForestGreen,
            urlcolor=MidnightBlue]{hyperref}

%% Headers / footers
\usepackage{fancyhdr}
\pagestyle{fancy}
\fancyhf{}
\rhead{\small ES-RL: Mathematical Derivation and Literature Review}
\lhead{\small \leftmark}
\cfoot{\thepage}
\renewcommand{\headrulewidth}{0.4pt}

%% Section styling
\usepackage{titlesec}
\titleformat{\section}{\large\bfseries\color{MidnightBlue}}{{\thesection.}}{0.5em}{}[\titlerule]
\titleformat{\subsection}{\normalsize\bfseries\color{MidnightBlue!80}}{{\thesubsection}}{0.5em}{}

%% Theorem environments
\newtheorem{theorem}{Theorem}[section]
\newtheorem{proposition}[theorem]{Proposition}
\newtheorem{lemma}[theorem]{Lemma}
\newtheorem{remark}{Remark}[section]

%% Coloured boxes
\newtcolorbox{keyresult}[1][]{
  enhanced, breakable,
  colback=MidnightBlue!6, colframe=MidnightBlue!60,
  fonttitle=\bfseries, title={Key Result}, #1
}
\newtcolorbox{importantbox}[1][]{
  enhanced, breakable,
  colback=ForestGreen!6, colframe=ForestGreen!60,
  fonttitle=\bfseries, title={#1}
}

%% ── Math macros (safe names — no clashes) ───────────────────
\newcommand{\bmu}{\boldsymbol{\mu}}
\newcommand{\btheta}{\boldsymbol{\theta}}
\newcommand{\beps}{\boldsymbol{\varepsilon}}
\newcommand{\bx}{\mathbf{x}}
\newcommand{\ghat}{\hat{\mathbf{g}}}     % gradient estimate
\newcommand{\mvec}{\mathbf{m}}           % first moment
\newcommand{\vvec}{\mathbf{v}}           % second moment
\newcommand{\mhat}{\hat{\mathbf{m}}}     % bias-corrected first moment
\newcommand{\vhat}{\hat{\mathbf{v}}}     % bias-corrected second moment
\newcommand{\EE}{\mathbb{E}}
\newcommand{\Var}{\mathrm{Var}}
\newcommand{\Cov}{\mathrm{Cov}}
\newcommand{\RR}{\mathbb{R}}
\newcommand{\NN}{\mathcal{N}}
\newcommand{\cX}{\mathcal{X}}
\DeclareMathOperator{\clip}{clip}

%% ── Document ────────────────────────────────────────────────
\begin{document}

%% ── Title page ──────────────────────────────────────────────
\begin{titlepage}
  \centering
  \vspace*{2cm}

  {\Huge\bfseries\color{MidnightBlue}
    Evolution Strategies\\[0.4em]
    with RL-style Gradient Estimation\\[0.4em]
    (ES-RL)
  }

  \vspace{0.8cm}
  \rule{0.7\linewidth}{1.5pt}

  \vspace{0.8cm}
  {\LARGE\itshape
    A Mathematical Derivation\\[0.3em]
    and Literature Review
  }

  \vspace{1.5cm}
  {\large\color{gray} Technical Document}

  \vspace{3cm}

  \begin{tcolorbox}[
    enhanced, colback=gray!8, colframe=MidnightBlue!40,
    width=0.85\linewidth, center, arc=4pt
  ]
    \textbf{Abstract.}
    This document presents a self-contained mathematical exposition of the
    Evolution Strategy with Reinforcement-Learning-style gradient estimation
    (ES-RL) algorithm.
    Starting from the first principles of stochastic optimisation, we derive
    the score-function (REINFORCE) gradient estimator, motivate antithetic
    sampling as a variance-reduction technique, and derive the Adam adaptive
    moment update.
    The complete ES-RL algorithm is assembled in a unified framework;
    connections to natural evolution strategies and policy-gradient methods
    are established; and convergence properties are summarised.
    All results are anchored to primary literature references.
  \end{tcolorbox}

  \vfill
  {\small Last compiled: \today}
\end{titlepage}

%% ── Table of Contents ───────────────────────────────────────
\tableofcontents
\newpage

% ================================================================
\section{Introduction and Motivation}
\label{sec:intro}
% ================================================================

Classical gradient-based optimisation requires access to $\nabla f(\bx)$,
either analytically or via automatic differentiation.
In many scientifically and industrially important settings this is not possible:

\begin{itemize}
  \item The objective $f$ is a \textbf{black box}
        (a physics simulation, an experiment, a quantum circuit).
  \item $f$ is \textbf{non-differentiable} or \textbf{discontinuous}.
  \item The search space involves \textbf{discrete decisions} mixed with
        continuous parameters.
  \item The computational graph is too large for backpropagation.
\end{itemize}

Evolution Strategies (ES) address all of these cases by estimating the
gradient of a \emph{smoothed} version of $f$ using only function evaluations.
The modern Reinforcement-Learning framing --- due primarily to
\textbf{Salimans et al.\ (2017)} \cite{salimans2017} and
\textbf{Wierstra et al.\ (2008, 2014)} \cite{wierstra2008, wierstra2014} ---
reinterprets ES as a policy-gradient algorithm on the distribution of
candidate solutions, enabling the use of powerful optimisers such as
\textbf{Adam} \cite{kingma2015} developed for deep learning.

% ================================================================
\section{Problem Formulation}
\label{sec:problem}
% ================================================================

We seek to \textbf{maximise} a scalar objective:
\begin{equation}
  \bx^{*} = \arg\max_{\bx \in \cX}\; f(\bx),
\end{equation}
where $f : \cX \subseteq \RR^d \to \RR$ may be \textbf{multimodal},
\textbf{noisy}, and \textbf{non-differentiable}.
Direct access to $\nabla f$ is \emph{not} assumed.

\subsection{Smoothing via a Search Distribution}

Rather than optimising $f$ directly, ES introduces a parametric family of
probability distributions $\{p_{\btheta}\}_{\btheta}$ over $\cX$ and
maximises the \textbf{expected fitness}:
\begin{equation}
  J(\btheta)
  = \EE_{\bx \sim p_{\btheta}}\!\bigl[f(\bx)\bigr]
  = \int_{\cX} f(\bx)\, p_{\btheta}(\bx)\, d\bx.
  \label{eq:expected_fitness}
\end{equation}
$J(\btheta)$ is a \textbf{smoothed, everywhere differentiable}
surrogate for $f$, even when $f$ itself is neither \cite{nesterov2017}.
Maximising $J(\btheta)$ concentrates probability mass near $\bx^*$.

\subsection{Isotropic Gaussian Search Distribution}

Throughout this document we use the \textbf{isotropic Gaussian}:
\begin{equation}
  p_{\btheta}(\bx) = \NN\!\left(\bx;\; \bmu,\; \sigma^2 \mathbf{I}\right),
\end{equation}
where $\btheta = \bmu \in \RR^d$ is the only optimisable parameter
(the \textbf{mean} of the search distribution) and
$\sigma > 0$ is a fixed \textbf{exploration spread}.
The goal reduces to:
\begin{quote}
  \itshape Find the $\bmu$ that maximises $J(\bmu)$.
\end{quote}

% ================================================================
\section{Evolution Strategies: Historical Background}
\label{sec:history}
% ================================================================

Evolution Strategies were introduced independently by
\textbf{Rechenberg} (1965, published 1973 \cite{rechenberg1973}) and
\textbf{Schwefel} (1965, published 1977 \cite{schwefel1977}) at the
Technical University of Berlin.
The original $(\mu,\lambda)$-ES:
\begin{enumerate}
  \item Generates $\lambda$ offspring from $\mu$ parents by adding Gaussian noise.
  \item Ranks offspring by fitness.
  \item Selects the $\rho$ best as the next generation's parents.
  \item Repeats.
\end{enumerate}
Step-size self-adaptation was formalised by \textbf{Hansen and Ostermeier}
in the Covariance Matrix Adaptation ES (CMA-ES) \cite{hansen2001}.
The modern \textbf{differentiable ES / Natural ES} framing
\cite{wierstra2008, wierstra2014} reinterprets the fitness-weighted
selection step as gradient ascent on $J(\btheta)$, enabling any
first-order optimiser.

% ================================================================
\section{The Score-Function Gradient Estimator}
\label{sec:score}
% ================================================================

\subsection{Derivation: The Log-Derivative Trick}

To perform gradient ascent on $J(\bmu)$ we need $\nabla_{\bmu} J$:
\begin{equation}
  \nabla_{\bmu} J(\bmu)
  = \nabla_{\bmu} \int f(\bx)\, p_{\bmu}(\bx)\, d\bx.
\end{equation}

Under mild regularity conditions (Lebesgue dominated convergence)
we may exchange the gradient and the integral:
\begin{equation}
  = \int f(\bx)\;\nabla_{\bmu} p_{\bmu}(\bx)\; d\bx.
\end{equation}

Applying the \textbf{log-derivative identity}
\begin{equation}
  \nabla_{\bmu} p_{\bmu}(\bx)
  = p_{\bmu}(\bx)\;\nabla_{\bmu} \log p_{\bmu}(\bx)
  \label{eq:logderiv}
\end{equation}
and substituting back, we obtain the \textbf{score-function estimator}:

\begin{keyresult}
  \begin{equation}
    \nabla_{\bmu} J(\bmu)
    = \EE_{\bx \sim p_{\bmu}}\!
      \left[f(\bx)\cdot \nabla_{\bmu} \log p_{\bmu}(\bx)\right].
    \label{eq:score}
  \end{equation}
\end{keyresult}

\noindent
This is also known as the \textbf{REINFORCE gradient} in reinforcement
learning \cite{williams1992} and the \textbf{natural gradient direction}
in natural evolution strategies \cite{wierstra2008, wierstra2014}.

\begin{remark}
  The right-hand side of \eqref{eq:score} is an expectation estimable by
  Monte Carlo sampling --- requiring \emph{only function evaluations},
  not analytic gradients of $f$.
\end{remark}

\subsection{Closed Form for the Isotropic Gaussian}

For $p_{\bmu}(\bx) = \NN(\bx;\bmu,\sigma^2\mathbf{I})$:
\begin{equation}
  \log p_{\bmu}(\bx)
  = -\tfrac{d}{2}\log(2\pi\sigma^2)
    - \frac{\|\bx - \bmu\|^2}{2\sigma^2}.
\end{equation}
\begin{equation}
  \nabla_{\bmu} \log p_{\bmu}(\bx)
  = \frac{\bx - \bmu}{\sigma^2}
  = \frac{\beps}{\sigma},
  \qquad
  \text{where } \bx = \bmu + \sigma\beps,\;
  \beps \sim \NN(\mathbf{0}, \mathbf{I}).
\end{equation}

Substituting into \eqref{eq:score}:
\begin{equation}
  \nabla_{\bmu} J(\bmu)
  = \frac{1}{\sigma}\;
    \EE_{\beps \sim \NN(\mathbf{0},\mathbf{I})}
    \!\left[f\!\left(\bmu + \sigma\beps\right)\cdot\beps\right].
  \label{eq:gaussian_grad}
\end{equation}

\subsection{Monte Carlo Estimator}

Using $H$ independent samples $\{\beps_i\}_{i=1}^{H}$:
\begin{equation}
  \hat{\mathbf{g}}_H
  = \frac{1}{H\sigma}\sum_{i=1}^{H}
    f\!\left(\bmu + \sigma\beps_i\right)\beps_i.
  \label{eq:mc}
\end{equation}

\begin{proposition}[Unbiasedness]
  $\EE\!\left[\hat{\mathbf{g}}_H\right] = \nabla_{\bmu} J(\bmu)$.
\end{proposition}

\noindent
The variance scales as:
\begin{equation}
  \Var\!\left[\hat{\mathbf{g}}_H\right]
  = \frac{1}{H\sigma^2}
    \Var_{\beps}\!\left[f(\bmu+\sigma\beps)\,\beps\right]
  = O\!\left(\tfrac{1}{H}\right),
\end{equation}
motivating the use of as many perturbations as the budget allows.

% ================================================================
\section{Antithetic Sampling: Variance Reduction}
\label{sec:antithetic}
% ================================================================

\subsection{Motivation}

The estimator \eqref{eq:mc} has variance $O(1/H)$.
For finite $H$, individual estimates can be very noisy when $f$ is
multimodal or when $\sigma$ is small relative to the scale of the landscape.

\subsection{Antithetic Pairs}

The \textbf{antithetic variates} method \cite{bratley1987} exploits the
symmetry of the Gaussian by evaluating \textbf{symmetric pairs}
$(\beps_i,-\beps_i)$:
\begin{equation}
  \hat{\mathbf{g}}_{\mathrm{anti}}
  = \frac{1}{2H\sigma}\sum_{i=1}^{H}
    \Bigl[f(\bmu+\sigma\beps_i) - f(\bmu-\sigma\beps_i)\Bigr]\,\beps_i.
  \label{eq:antithetic}
\end{equation}

\subsection{Proof of Unbiasedness}

\begin{proof}
Since $\beps_i \overset{d}{=} -\beps_i$ (Gaussian symmetry):
\begin{align}
  \EE\!\left[\hat{\mathbf{g}}_{\mathrm{anti}}\right]
  &= \frac{1}{2\sigma}
     \EE\!\left[f(\bmu+\sigma\beps)\beps\right]
   + \frac{1}{2\sigma}
     \EE\!\left[-f(\bmu-\sigma\beps)\beps\right]
  \nonumber\\
  &= \frac{1}{2\sigma}
     \EE\!\left[f(\bmu+\sigma\beps)\beps\right]
   + \frac{1}{2\sigma}
     \EE\!\left[f(\bmu+\sigma\beps)\beps\right]
  \nonumber\\
  &= \frac{1}{\sigma}
     \EE\!\left[f(\bmu+\sigma\beps)\beps\right]
   = \nabla_{\bmu} J(\bmu). \qquad\checkmark
\end{align}
The second equality uses the substitution $\beps \mapsto -\beps$,
which leaves the Gaussian distribution invariant.
\end{proof}

\subsection{Variance Reduction Analysis}

Let $A_i = f(\bmu+\sigma\beps_i)\varepsilon_i$ and
$B_i = -f(\bmu-\sigma\beps_i)\varepsilon_i$.
The antithetic estimate uses $(A_i+B_i)/2$ instead of $A_i$.
For smooth $f$, a positive perturbation that yields a high value makes
the mirror perturbation tend to yield a lower value; hence
$\Cov(A_i,B_i)<0$, and:
\begin{equation}
  \Var\!\left[\frac{A_i+B_i}{2}\right]
  = \frac{\Var[A_i]+\Var[B_i]+2\Cov(A_i,B_i)}{4}
  < \frac{\Var[A_i]}{2}.
\end{equation}

\begin{keyresult}
  Antithetic sampling \textbf{at least halves the estimator variance}
  relative to one-sided sampling, at no extra cost in the number of
  perturbation \emph{directions} $H$ (though $2H$ total evaluations
  are required).
\end{keyresult}

\noindent
This was demonstrated empirically by Salimans et al.\ \cite{salimans2017},
who showed that antithetic sampling roughly doubles effective sample
efficiency in high-dimensional problems.

% ================================================================
\section{The Adam Optimiser}
\label{sec:adam}
% ================================================================

\subsection{Motivation}

Applying the raw estimate $\hat{\mathbf{g}}$ in a fixed-step update
$\bmu \leftarrow \bmu + \alpha\hat{\mathbf{g}}$ suffers from:
\begin{itemize}
  \item \textbf{Sensitivity to $\alpha$:}
        too large $\Rightarrow$ oscillation; too small $\Rightarrow$ slow convergence.
  \item \textbf{Non-stationarity:}
        the optimal step size changes as $\bmu$ approaches the optimum.
  \item \textbf{Noise amplification:}
        estimates from a small batch of perturbations are inherently noisy.
\end{itemize}
\textbf{Adam} (Adaptive Moment Estimation) \cite{kingma2015} addresses all
three by maintaining per-parameter running statistics of the gradient.

\subsection{Derivation from Exponential Moving Averages}

Let $\hat{\mathbf{g}}_t$ be the gradient estimate at step $t$.  Define:

\medskip
\noindent
\textbf{First moment} (momentum):
\begin{equation}
  \mvec_t = \beta_1\,\mvec_{t-1} + (1-\beta_1)\,\hat{\mathbf{g}}_t,
  \qquad \mvec_0 = \mathbf{0}.
  \label{eq:m}
\end{equation}

\noindent
\textbf{Second moment} (uncentred variance):
\begin{equation}
  \vvec_t = \beta_2\,\vvec_{t-1} + (1-\beta_2)\,\hat{\mathbf{g}}_t^{\,2},
  \qquad \vvec_0 = \mathbf{0},
  \label{eq:v}
\end{equation}
where $\beta_1,\beta_2\in[0,1)$ and squaring is elementwise.

\subsection{Bias Correction}

Unrolling \eqref{eq:m} from $\mvec_0=\mathbf{0}$:
\begin{equation}
  \mvec_t = (1-\beta_1)\sum_{i=1}^{t}\beta_1^{\,t-i}\,\hat{\mathbf{g}}_i.
\end{equation}
Taking expectations under the stationarity assumption
$\EE[\hat{\mathbf{g}}_i]=\mathbf{g}$:
\begin{equation}
  \EE[\mvec_t]
  = \mathbf{g}\,(1-\beta_1)\sum_{i=1}^{t}\beta_1^{\,t-i}
  = \mathbf{g}\,(1-\beta_1^t).
\end{equation}
The estimate is \textbf{biased toward zero} by the factor $(1-\beta_1^t)$.
Analogously, $\EE[\vvec_t] = \EE[\mathbf{g}^2]\,(1-\beta_2^t)$.
We correct with:
\begin{equation}
  \mhat_t = \frac{\mvec_t}{1-\beta_1^t},
  \qquad
  \vhat_t = \frac{\vvec_t}{1-\beta_2^t},
\end{equation}
so that $\EE[\mhat_t]\approx\mathbf{g}$ and
$\EE[\vhat_t]\approx\EE[\mathbf{g}^2]$.

\subsection{The Update Rule}

\begin{keyresult}
  \begin{equation}
    \bmu_{t+1}
    = \bmu_t
      + \alpha\cdot
        \frac{\mhat_t}{\sqrt{\vhat_t}+\varepsilon},
    \label{eq:adam}
  \end{equation}
  where $\varepsilon\approx10^{-8}$ prevents division by zero.
\end{keyresult}

\begin{table}[ht]
  \centering
  \caption{Interpretation of the Adam update terms.}
  \label{tab:adam_terms}
  \begin{tabular}{@{}lp{9.5cm}@{}}
    \toprule
    \textbf{Term} & \textbf{Role} \\
    \midrule
    $\mhat_t$ &
      Smoothed gradient direction; reduces noise via momentum, suppresses
      short-lived fluctuations. \\[4pt]
    $\sqrt{\vhat_t}$ &
      Running RMS of gradient magnitude; normalises the step size
      per coordinate. \\[4pt]
    $\alpha/\sqrt{\vhat_t}$ &
      Per-parameter adaptive rate: smaller for high-gradient coordinates,
      larger for low-gradient or infrequently-updated ones. \\
    \bottomrule
  \end{tabular}
\end{table}

\subsection{Standard Hyperparameters}

\begin{table}[ht]
  \centering
  \caption{Recommended Adam hyperparameters \cite{kingma2015}.}
  \label{tab:adam_hyp}
  \begin{tabular}{@{}ccc@{}}
    \toprule
    \textbf{Parameter} & \textbf{Typical value} & \textbf{Role} \\
    \midrule
    $\alpha$      & $10^{-3}$--$0.05$ & Global learning rate (primary tuning knob) \\
    $\beta_1$     & $0.9$             & Momentum decay \\
    $\beta_2$     & $0.999$           & RMS decay \\
    $\varepsilon$ & $10^{-8}$         & Numerical stability \\
    \bottomrule
  \end{tabular}
\end{table}

% ================================================================
\section{The Complete ES-RL Algorithm}
\label{sec:algorithm}
% ================================================================

\subsection{State Variables}

\begin{table}[ht]
  \centering
  \caption{ES-RL state variables.}
  \begin{tabular}{@{}lll@{}}
    \toprule
    \textbf{Variable} & \textbf{Dimension} & \textbf{Description} \\
    \midrule
    $\bmu$       & $d$     & Current mean of the search distribution \\
    $\sigma$     & scalar  & Exploration spread (fixed) \\
    $\mvec$      & $d$     & Adam first-moment accumulator \\
    $\vvec$      & $d$     & Adam second-moment accumulator \\
    $t$          & integer & Adam timestep counter \\
    \bottomrule
  \end{tabular}
\end{table}

\subsection{One Iteration}

\noindent\textbf{Step 1 — Sample perturbations:}
\begin{equation*}
  \beps_1,\ldots,\beps_H \overset{i.i.d.}{\sim} \NN(\mathbf{0},\mathbf{I}_d).
\end{equation*}

\noindent\textbf{Step 2 — Evaluate antithetic pairs} ($2H$ evaluations):
\begin{equation*}
  f^+_i = f(\bmu+\sigma\beps_i),\quad
  f^-_i = f(\bmu-\sigma\beps_i),\qquad i=1,\ldots,H.
\end{equation*}

\noindent\textbf{Step 3 — Gradient estimate} (eq.~\eqref{eq:antithetic}):
\begin{equation*}
  \hat{\mathbf{g}}
  = \frac{1}{2H\sigma}\sum_{i=1}^{H}(f^+_i-f^-_i)\,\beps_i.
\end{equation*}

\noindent\textbf{Step 4 — Adam moment update:}
\begin{equation*}
  t\leftarrow t+1,\quad
  \mvec\leftarrow\beta_1\mvec+(1-\beta_1)\hat{\mathbf{g}},\quad
  \vvec\leftarrow\beta_2\vvec+(1-\beta_2)\hat{\mathbf{g}}^{\,2}.
\end{equation*}

\noindent\textbf{Step 5 — Bias correction:}
\begin{equation*}
  \mhat = \mvec/(1-\beta_1^t),\qquad
  \vhat = \vvec/(1-\beta_2^t).
\end{equation*}

\noindent\textbf{Step 6 — Update mean:}
\begin{equation*}
  \bmu\leftarrow
  \clip\!\left(
    \bmu+\alpha\cdot\frac{\mhat}{\sqrt{\vhat}+\varepsilon},\;
    \bx_{\min},\;\bx_{\max}
  \right).
\end{equation*}

\subsection{Full Pseudocode}

\begin{algorithm}[H]
  \caption{ES-RL with Antithetic Sampling and Adam}
  \label{alg:esrl}
  \begin{algorithmic}[1]
    \Require $f,\;\cX{=}[\bx_{\min},\bx_{\max}],\;
             \sigma,\;H,\;\alpha,\;
             \beta_1{=}0.9,\;\beta_2{=}0.999,\;
             \varepsilon{=}10^{-8}$, budget $N$
    \State $\bmu \leftarrow \text{random point in } \cX$
    \State $\mvec\leftarrow\mathbf{0};\;
           \vvec\leftarrow\mathbf{0};\;
           t\leftarrow 0$
           \Comment{Adam state}
    \State $f^*\leftarrow-\infty;\;\bx^*\leftarrow\text{None}$
    \While{evaluations $< N$}
      \State Sample $\beps_1,\ldots,\beps_H \overset{i.i.d.}{\sim}
             \NN(\mathbf{0},\mathbf{I})$
      \For{$i=1,\ldots,H$}
        \State $f^+_i\leftarrow f\!\left(\clip(\bmu+\sigma\beps_i,
               \bx_{\min},\bx_{\max})\right)$
        \State $f^-_i\leftarrow f\!\left(\clip(\bmu-\sigma\beps_i,
               \bx_{\min},\bx_{\max})\right)$
        \State Update $f^*,\bx^*$ if $\max(f^+_i,f^-_i)>f^*$
      \EndFor
      \State $\hat{\mathbf{g}}\leftarrow
             \dfrac{1}{2H\sigma}\displaystyle\sum_{i=1}^{H}
             (f^+_i-f^-_i)\,\beps_i$
      \State $t\leftarrow t+1$
      \State $\mvec\leftarrow\beta_1\mvec+(1-\beta_1)\hat{\mathbf{g}};\;
             \vvec\leftarrow\beta_2\vvec+(1-\beta_2)\hat{\mathbf{g}}^2$
      \State $\mhat\leftarrow\mvec/(1-\beta_1^t);\;
             \vhat\leftarrow\vvec/(1-\beta_2^t)$
      \State $\bmu\leftarrow\clip\!\bigl(
             \bmu+\alpha\,\mhat/(\sqrt{\vhat}+\varepsilon),\,
             \bx_{\min},\bx_{\max}\bigr)$
    \EndWhile
    \State\Return $\bx^*,f^*$
  \end{algorithmic}
\end{algorithm}

% ================================================================
\section{Convergence Analysis}
\label{sec:convergence}
% ================================================================

\subsection{Relationship to Gaussian Smoothing}

The expected fitness $J(\bmu)$ equals the Gaussian convolution of $f$:
\begin{equation}
  J(\bmu) = (f*\phi_\sigma)(\bmu),\qquad
  \phi_\sigma(\bx) = \NN(\bx;\,\mathbf{0},\sigma^2\mathbf{I}).
\end{equation}
Nesterov and Spokoiny \cite{nesterov2017} showed that $J$ inherits Lipschitz
smoothness from $f$: if $f$ is $L$-Lipschitz, then $J$ is $L$-Lipschitz
and its gradient is $(L/\sigma)$-Lipschitz, guaranteeing convergence
of gradient ascent on $J$.

\subsection{Strongly Concave Case}

For $J$ that is $\kappa$-strongly concave near the maximum, projected
gradient ascent with step $\alpha$ converges geometrically \cite{bubeck2015}:
\begin{equation}
  J(\bmu^*)-J(\bmu_t)
  \;\leq\;
  (1-\alpha\kappa)^t\,\bigl[J(\bmu^*)-J(\bmu_0)\bigr].
\end{equation}
In practice $J$ inherits the multimodality of $f$ and convergence to a
\emph{local} maximum is what is guaranteed.

\subsection{Finite-Sample Rate (Non-Convex)}

With $H$ antithetic pairs, the gradient estimate has variance $O(1/H)$.
By the non-convex stochastic gradient theory \cite{bottou2018}:
\begin{equation}
  \frac{1}{T}\sum_{t=1}^{T}
  \EE\!\left[\bigl\|\nabla J(\bmu_t)\bigr\|^2\right]
  = O\!\left(\frac{1}{\sqrt{T}}+\frac{\sigma_g}{\sqrt{H}}\right),
\end{equation}
where $\sigma_g$ is the variance of one gradient sample.
Larger $H$ tightens the bound.

\subsection{Approximation Bias from $\sigma$}

A Taylor expansion \cite{salimans2017} shows:
\begin{equation}
  \nabla_{\bmu}J(\bmu) = \nabla f(\bmu) + O(\sigma^2).
\end{equation}
Small $\sigma$: nearly unbiased, but high variance and inability to
escape local maxima.
Large $\sigma$: low variance but the smoothed surrogate may miss narrow
peaks.
This tension is the fundamental exploration--exploitation trade-off in ES.

% ================================================================
\section{Connections to Related Methods}
\label{sec:connections}
% ================================================================

\subsection{REINFORCE (Williams, 1992)}

Williams \cite{williams1992} derived the score-function estimator in
reinforcement learning, where $f(\bx)$ is a \textbf{return} and
$p_{\btheta}(\bx)$ is a \textbf{stochastic policy}:
\begin{equation}
  \hat{\mathbf{g}}_{\mathrm{REINFORCE}}
  = \frac{1}{H}\sum_{i=1}^{H}
    R(\tau_i)\,\nabla_{\btheta}\log\pi_{\btheta}(\tau_i).
\end{equation}
Substituting $R\equiv f$, $\pi_{\btheta}\equiv p_{\bmu}$,
$\tau_i\equiv\bx_i$ recovers eq.~\eqref{eq:mc} exactly.
ES-RL is a direct instantiation of REINFORCE with a Gaussian ``policy.''

\subsection{Natural Evolution Strategies (NES)}

Wierstra et al.\ \cite{wierstra2008, wierstra2014} proposed the
\textbf{natural gradient} update:
\begin{equation}
  \bmu_{t+1} = \bmu_t + \alpha\,\mathbf{F}^{-1}(\bmu_t)\,\hat{\mathbf{g}},
\end{equation}
where $\mathbf{F}(\bmu)=\EE\!\bigl[(\nabla_{\bmu}\log p_{\bmu})
(\nabla_{\bmu}\log p_{\bmu})^\top\bigr]$ is the Fisher information matrix.
For the isotropic Gaussian $\mathbf{F}=\sigma^{-2}\mathbf{I}$, which
simplifies to $\bmu_{t+1}=\bmu_t+(\alpha/\sigma)\hat{\mathbf{g}}_{\mathrm{norm}}$,
removing the scale dependence on $\sigma$.

\subsection{CMA-ES}

CMA-ES \cite{hansen2001} adapts a full covariance matrix
$\mathbf{C}$ giving second-order-like convergence at $O(d^2)$ cost per step.
ES-RL with an isotropic Gaussian is a first-order $O(d)$ variant, trading
convergence speed for scalability.

% ================================================================
\section{Summary and Practical Guidance}
\label{sec:summary}
% ================================================================

\subsection{Algorithm in One Equation}

\begin{importantbox}[ES-RL Update (complete)]
  \begin{equation}
    \bmu_{t+1}
    = \clip\!\left(
        \bmu_t
        + \alpha\cdot
          \frac{\mhat_t}{\sqrt{\vhat_t}+\varepsilon},\;
        \bx_{\min},\;\bx_{\max}
      \right),
    \label{eq:full_update}
  \end{equation}
  where $\mhat_t$ and $\vhat_t$ are the bias-corrected Adam moments of
  the antithetic gradient estimate \eqref{eq:antithetic}.
\end{importantbox}

\subsection{Hyperparameter Guidance}

\begin{table}[ht]
  \centering
  \caption{ES-RL hyperparameter guidance.}
  \label{tab:hyp}
  \begin{tabular}{@{}p{3.5cm}p{3cm}p{7cm}@{}}
    \toprule
    \textbf{Hyperparameter} & \textbf{Recommended} & \textbf{Effect of increasing} \\
    \midrule
    $\sigma$ (spread) & 0.05--0.20 (fraction of domain) &
      Better exploration; more approximation bias \\[3pt]
    $H$ (pairs/step)  & 5--20 &
      Lower gradient variance; more evaluations per step \\[3pt]
    $\alpha$ (learning rate) & 0.01--0.10 &
      Faster but potentially less stable convergence \\[3pt]
    $\beta_1$ (momentum) & 0.9 (fixed) &
      Smoother gradient direction \\[3pt]
    $\beta_2$ (RMS decay) & 0.999 (fixed) &
      More stable norm estimation \\
    \bottomrule
  \end{tabular}
\end{table}

\subsection{When ES-RL Outperforms Random Search}

\textbf{ES-RL is advantageous when:}
\begin{itemize}
  \item The landscape has a \textbf{clear basin} around the optimum
        (gradient signal is informative).
  \item The \textbf{evaluation budget is limited} relative to domain size.
  \item Dimensionality $d>1$ (random search needs
        $O((1/p_\varepsilon)^{1/d})$ samples --- exponential in $d$).
\end{itemize}

\textbf{Random search remains competitive when:}
\begin{itemize}
  \item The landscape is \textbf{pathological} (needle-in-haystack;
        gradient information is useless).
  \item The domain is small or \textbf{1-dimensional} with modest budget.
  \item \textbf{Global diversity} is more important than local convergence.
\end{itemize}

% ================================================================
\section*{References}
\addcontentsline{toc}{section}{References}
% ================================================================

\begin{thebibliography}{15}

\bibitem{salimans2017}
T.~Salimans, J.~Ho, X.~Chen, S.~Sidor, and I.~Sutskever,
\newblock ``Evolution Strategies as a Scalable Alternative to Reinforcement
  Learning,''
\newblock \textit{arXiv:1703.03864}, OpenAI, 2017.

\bibitem{wierstra2008}
D.~Wierstra, T.~Forster, J.~Peters, and J.~Schmidhuber,
\newblock ``Natural Evolution Strategies,''
\newblock \textit{Proc.\ IEEE Congress on Evolutionary Computation (CEC)},
  2008.

\bibitem{wierstra2014}
D.~Wierstra, T.~Forster, J.~Peters, and J.~Schmidhuber,
\newblock ``Natural Evolution Strategies,''
\newblock \textit{Journal of Machine Learning Research},
  \textbf{15}: 949--980, 2014.

\bibitem{kingma2015}
D.~P.~Kingma and J.~Ba,
\newblock ``Adam: A Method for Stochastic Optimization,''
\newblock \textit{Proc.\ ICLR}, 2015.

\bibitem{nesterov2017}
Y.~Nesterov and V.~Spokoiny,
\newblock ``Random Gradient-Free Minimization of Convex Functions,''
\newblock \textit{Foundations of Computational Mathematics},
  \textbf{17}(2): 527--566, 2017.

\bibitem{rechenberg1973}
I.~Rechenberg,
\newblock \textit{Evolutionsstrategie: Optimierung technischer Systeme nach
  Prinzipien der biologischen Evolution},
\newblock Friedrich Frommann Verlag, Stuttgart, 1973.

\bibitem{schwefel1977}
H.-P.~Schwefel,
\newblock \textit{Numerische Optimierung von Computer-Modellen mittels der
  Evolutionsstrategie},
\newblock Birkh\"{a}user, Basel, 1977.

\bibitem{hansen2001}
N.~Hansen and A.~Ostermeier,
\newblock ``Completely Derandomized Self-Adaptation in Evolution Strategies,''
\newblock \textit{Evolutionary Computation}, \textbf{9}(2): 159--195, 2001.

\bibitem{williams1992}
R.~J.~Williams,
\newblock ``Simple Statistical Gradient-Following Algorithms for
  Connectionist Reinforcement Learning,''
\newblock \textit{Machine Learning}, \textbf{8}(3--4): 229--256, 1992.

\bibitem{bratley1987}
P.~Bratley, B.~L.~Fox, and L.~E.~Schrage,
\newblock \textit{A Guide to Simulation}, 2nd~ed.
\newblock Springer, New York, 1987.

\bibitem{bubeck2015}
S.~Bubeck,
\newblock ``Convex Optimization: Algorithms and Complexity,''
\newblock \textit{Foundations and Trends in Machine Learning},
  \textbf{8}(3--4): 231--357, 2015.

\bibitem{bottou2018}
L.~Bottou, F.~E.~Curtis, and J.~Nocedal,
\newblock ``Optimization Methods for Large-Scale Machine Learning,''
\newblock \textit{SIAM Review}, \textbf{60}(2): 223--311, 2018.

\bibitem{sutton2018}
R.~S.~Sutton and A.~G.~Barto,
\newblock \textit{Reinforcement Learning: An Introduction}, 2nd~ed.,
\newblock MIT Press, 2018, Ch.~13.

\bibitem{hansen2016}
N.~Hansen,
\newblock ``The CMA Evolution Strategy: A Tutorial,''
\newblock \textit{arXiv:1604.00772}, 2016.

\bibitem{graves2013}
A.~Graves,
\newblock ``Generating Sequences with Recurrent Neural Networks,''
\newblock \textit{arXiv:1308.0850}, 2013.

\end{thebibliography}

\end{document}
